% ESP8266 (NodeMCU V3 / LoLin) Relay & Sensor Controller -- schematic, Rev C
% Build: pdflatex schematic.tex   (needs the circuitikz package)
\documentclass[border=12pt]{standalone}
\usepackage[american]{circuitikz}
\usepackage{xcolor}

\definecolor{rail5}{RGB}{200,30,30}    % +5V net
\definecolor{rail3}{RGB}{230,120,0}    % +3V3 net
\definecolor{railg}{RGB}{40,40,40}     % GND net
\definecolor{sig}{RGB}{20,80,170}      % GPIO signals
\definecolor{nc}{RGB}{150,150,150}     % not connected

\ctikzset{bipoles/length=1.1cm}
\tikzset{
  mod/.style   = {draw, thick, fill=gray!8, rounded corners=2pt},
  pin/.style   = {font=\scriptsize\sffamily},
  title/.style = {font=\small\sffamily\bfseries},
  note/.style  = {font=\scriptsize\sffamily, align=left},
  p5/.style    = {rail5, very thick},
  p3/.style    = {rail3, very thick},
  pg/.style    = {railg, very thick},
  s/.style     = {sig, thick},
}

% ---- U1 pin helpers (box spans x = 0 .. 5) --------------------------------
% \lpin{y}{label}, \rpin{y}{label}: connected pin
% \lnc{y}{label},  \rnc{y}{label} : pin present on the board, not connected
\newcommand{\lpin}[2]{\draw[thick] (-0.6,#1) -- (0,#1);
  \node[pin, anchor=west] at (0.1,#1) {#2};}
\newcommand{\rpin}[2]{\draw[thick] (5,#1) -- (5.6,#1);
  \node[pin, anchor=east] at (4.9,#1) {#2};}
\newcommand{\xmark}[2]{\draw[nc, thick] (#1-0.1,#2-0.1) -- (#1+0.1,#2+0.1)
  (#1-0.1,#2+0.1) -- (#1+0.1,#2-0.1);}
\newcommand{\lnc}[2]{\draw[nc] (-0.5,#1) -- (0,#1); \xmark{-0.5}{#1}
  \node[pin, nc, anchor=west] at (0.1,#1) {#2};}
\newcommand{\rnc}[2]{\draw[nc] (5,#1) -- (5.5,#1); \xmark{5.5}{#1}
  \node[pin, nc, anchor=east] at (4.9,#1) {#2};}

% Relay module: \relay{IN y}{index}{MCU pin y}{bend x}{pin name}
\newcommand{\relay}[5]{
  \draw[mod] (11,#1-0.6) rectangle (15.4,#1+0.6);
  \node[title] at (13.2,#1+0.2) {K#2 \; RELAY #2 (5 V)};
  \node[pin]   at (13.2,#1-0.28) {COM / NO / NC $\to$ Load #2};
  \draw[s] (5.6,#3) -| (#4,#1) -- (11,#1);
  \node[pin, sig, anchor=south] at (10.2,#1) {#5 $\to$ IN};
  \draw[p5] (15.4,#1+0.3) -- (19,#1+0.3) node[circ]{};
  \node[pin, anchor=south west] at (15.45,#1+0.3) {VCC};
  \draw[pg] (15.4,#1-0.3) -- (16.8,#1-0.3) node[ground]{};
  \node[pin, anchor=north west] at (15.45,#1-0.3) {GND};
}

\begin{document}
\begin{circuitikz}

%======================================================================
% U1 -- ESP8266 NodeMCU V3 (LoLin), all 30 pins in board order
% (antenna at top, micro-USB at bottom)
%======================================================================
\draw[mod, fill=blue!4] (0,1.2) rectangle (5,16.2);
\draw[gray, thin] (1.4,15.0) rectangle (3.6,16.0);
\node[pin, gray] at (2.5,15.5) {ESP-12E (antenna)};
\node[title, anchor=south] at (2.5,16.3) {U1 \; ESP8266 NodeMCU V3 (LoLin)};

% Left header (top -> bottom)
\lpin{14.0}{A0 / ADC0}
\lnc{13.2}{G}
\lnc{12.4}{VU (VUSB)}
\lnc{11.6}{S3 / GPIO10}
\lnc{10.8}{S2 / GPIO9}
\lnc{10.0}{S1 / MOSI}
\lnc{9.2}{SC / CS}
\lnc{8.4}{S0 / MISO}
\lnc{7.6}{SK / SCLK}
\lnc{6.8}{G}
\lpin{6.0}{3V3}
\lnc{5.2}{EN}
\lnc{4.4}{RST}
\lnc{3.6}{G}
\lpin{2.8}{VIN}

% Right header (top -> bottom)
\rpin{14.0}{D0 / GPIO16}
\rpin{13.2}{D1 / GPIO5}
\rpin{12.4}{D2 / GPIO4}
\rpin{11.6}{D3 / GPIO0}
\rpin{10.8}{D4 / GPIO2}
\rnc{10.0}{3V3}
\rnc{9.2}{G}
\rpin{8.4}{D5 / GPIO14}
\rpin{7.6}{D6 / GPIO12}
\rpin{6.8}{D7 / GPIO13}
\rnc{6.0}{D8 / GPIO15}
\rnc{5.2}{RX / GPIO3}
\rnc{4.4}{TX / GPIO1}
\rpin{3.6}{G}
\rnc{2.8}{3V3}

% J1 micro-USB (bottom of board)
\draw[thick, fill=gray!25] (1.1,1.3) rectangle (3.9,2.2);
\node[pin] at (2.5,1.75) {J1 micro-USB (5 V)};

%======================================================================
% POWER RAILS
%======================================================================
% +5V: VIN -> right breadboard + rail
\draw[p5] (-0.6,2.8) -- (-4.3,2.8) arc(0:180:0.2) -- (-6,2.8) -- (-6,20) -- (19,20) -- (19,14.2);
\node[pin, rail5, anchor=south west] at (-5.9,20) {VIN = +5V};
\node[title, rail5, anchor=south] at (19,20.1) {+5V RAIL (breadboard right +)};

% +3V3: U1 3V -> left breadboard + rail
\draw[p3] (-0.6,6) -- (-4.5,6) -- (-4.5,-0.4) -- (18,-0.4) -- (18,11.9);
\node[pin, rail3, anchor=south east] at (-0.7,6) {+3V3 out};
\node[title, rail3, anchor=north] at (6.5,-0.5) {+3V3 from U1 left-header 3V3 (pin 11) $\to$ sensor rail};
\node[title, rail3, anchor=west, align=left] at (18.2,9) {+3V3 RAIL\\(breadboard left +)};

% Common GND bus
\draw[pg] (5.6,3.6) -- (6.2,3.6) -- (6.2,0.4);
\draw[pg] (6.2,0.4) -- (17,0.4) node[ground]{};
\node[title, railg, anchor=south west] at (7,0.45)
  {COMMON GND BUS -- both breadboard $-$ rails (every $\perp$ = this net)};

% C1 bulk filter across +5V
\draw[p5] (19,19.4) node[circ]{} -- (20.8,19.4);
\draw (20.8,19.4) to[eC, l_={C1}, a^={1000\,$\mu$F / 16\,V}, *-] (20.8,17.6) node[ground]{};

%======================================================================
% VIN MONITOR DIVIDER (A0)
%======================================================================
\draw[s] (-0.6,14) -- (-3,14) node[circ]{};
\draw (-3,20) node[circ, rail5]{} -- (-3,17.4) to[R, l_={R1}, a^={47\,k$\Omega$}] (-3,15.6) -- (-3,14);
\draw (-3,14) to[R, l_={R2}, a^={20\,k$\Omega$}] (-3,12.2) node[ground]{};
\node[title, anchor=east] at (-3.3,18.8) {VIN MONITOR};

%======================================================================
% RELAY BLOCK (5 V)
%======================================================================
\node[title, rail5] at (13.2,19.45) {RELAY BLOCK (5 V)};
\relay{18.4}{1}{14.0}{7.0}{D0}
\relay{16.9}{2}{13.2}{7.8}{D1}
\relay{15.4}{3}{12.4}{8.6}{D2}
\relay{13.9}{4}{11.6}{9.4}{D3}

%======================================================================
% SENSOR BLOCK (3.3 V)
%======================================================================
\node[title, rail3] at (13,12.75) {SENSOR BLOCK (3.3 V)};

% U2 DHT11
\draw[mod] (11,10.5) rectangle (15,12.3);
\node[title] at (13,11.8) {U2 \; DHT11};
\node[pin]   at (13,11.25) {Temp / Humidity};
\draw[s] (5.6,10.8) -- (11,10.8);
\node[pin, sig, anchor=south] at (10.2,10.8) {DATA};
\draw[p3] (15,11.9) -- (18,11.9);
\node[pin, anchor=south west] at (15.05,11.9) {VCC};
\draw[pg] (15,11.0) -- (16.3,11.0) node[ground]{};
\node[pin, anchor=north west] at (15.05,11.0) {GND};

% U3 HR202 module (LM393 comparator board)
\draw[mod] (13,5.6) rectangle (16.6,7.4);
\node[title] at (14.8,6.9) {U3 \; HR202};
\node[pin]   at (14.8,6.35) {Humidity module};
\draw[s] (5.6,6.8) -- (13,6.8);
\node[pin, sig, anchor=south] at (12.4,6.8) {DO};
\draw[nc] (13,6.0) -- (12.5,6.0); \xmark{12.5}{6.0}
\node[pin, nc, anchor=east] at (12.35,6.0) {AO};
\draw[p3] (16.6,7.1) -- (18,7.1);
\node[pin, anchor=south west] at (16.65,7.1) {VCC};
\draw[pg] (16.6,6.0) -- (17.6,6.0) node[ground]{};
\node[pin, anchor=north west] at (16.65,6.0) {GND};

%======================================================================
% INPUT BUTTONS (active low, INPUT_PULLUP)
%======================================================================
\node[title] at (9,9.35) {INPUTS (active low, INPUT\_PULLUP)};
\draw[s] (5.6,8.4) -- (6.4,8.4);
\draw (6.4,8.4) to[push button, l={SW1 -- Relay 1 ON}, *-] (8.2,8.4) -- (8.6,8.4) node[ground]{};
\draw[s] (5.6,7.6) -- (10,7.6);
\draw (10,7.6) to[push button, l={SW2 -- ALL OFF}, *-] (11.8,7.6) -- (12.2,7.6) node[ground]{};

%======================================================================
% LEGEND / TITLE BLOCK
%======================================================================
\draw[thick] (-7,-1.6) rectangle (22.4,-3.6);
\node[note, anchor=north west] at (-6.8,-1.75) {%
  \textbf{ESP8266 Relay \& Sensor Controller} \quad Rev C \quad (U1 = NodeMCU V3 / LoLin, all 30 header pins shown)\\[2pt]
  \textcolor{rail5}{\rule[2pt]{0.8cm}{1.5pt}} +5V (VIN) \quad
  \textcolor{rail3}{\rule[2pt]{0.8cm}{1.5pt}} +3V3 \quad
  \textcolor{railg}{\rule[2pt]{0.8cm}{1.5pt}} GND \quad
  \textcolor{sig}{\rule[2pt]{0.8cm}{1.5pt}} GPIO signal \quad
  \textcolor{nc}{$\times$} = pin not connected\\[2pt]
  Dot = junction; crossing without dot = no connection. All G pins are common internally; all 3V3 pins are common internally. Sensors use the LEFT 3V3 pin.\\
  Divider: 5 V $\to$ 1.49 V at A0 (NodeMCU A0 accepts 0--3.2 V via its onboard 220k/100k divider).\\
  Boot straps: D3 (GPIO0), D4 (GPIO2) must be HIGH and D8 (GPIO15) LOW at reset. RX/TX are used by the USB-serial chip.};

\end{circuitikz}
\end{document}
